\expandafter\def\csname HistoricalNoteText18\endcsname{Da dieser Wert von der zwanzigsten Ziffer an von demjenigen abweicht, den der berühmte \textsc{Mascheroni} in Adnotat. ad \textsc{Euleri} \textit{Calculum Integr.} berechnet hat, habe ich \textsc{Fridericus Bernhardus Gothofredus Nicolai}, einen im Rechnen unermüdlichen jungen Mann, dazu aufgefordert, jene Berechnung zu wiederholen und weiter fortzuführen. Er fand also durch zwei Berechnungen, nämlich sowohl von $z=50$ als auch von $z=100$ abwärts rechnend,
\[
\Psi0=-0{,}5772156649\ 0153286060\ 6512090082\ 4024310421
\]
Auch der zweite Teil der am Ende dieses Abschnitts angefügten Tafel stammt von demselben höchst geübten Rechner; er gibt die Werte der Funktion $\Psi z$ auf 18 Ziffern (deren letzte nicht sicher ist) für alle Werte von $z$ von 0 bis 1 in Schritten von je einem Hundertstel an. Im Übrigen beruhen die Methoden, nach denen beide Tafeln berechnet wurden, teils auf den hier mitgeteilten Sätzen, teils auf besonderen Kunstgriffen des Rechnens, die wir bei anderer Gelegenheit vorlegen werden.}
